\documentclass[11pt,a4paper]{article}

\usepackage[margin=2.2cm]{geometry}
\usepackage{kotex}
\usepackage{amsmath,amssymb}
\usepackage{booktabs}
\usepackage{array}
\usepackage{xcolor}
\usepackage{listings}
\usepackage[most]{tcolorbox}
\usepackage{enumitem}
\usepackage{hyperref}
\usepackage{fancyhdr}
\usepackage{titlesec}

\definecolor{codebg}{RGB}{247,247,247}
\definecolor{codeframe}{RGB}{210,210,210}
\definecolor{keywordcolor}{RGB}{0,70,160}
\definecolor{commentcolor}{RGB}{0,120,80}
\definecolor{stringcolor}{RGB}{170,50,50}
\definecolor{titleblue}{RGB}{35,75,130}

\hypersetup{colorlinks=true,linkcolor=titleblue,urlcolor=titleblue}

\pagestyle{fancy}
\fancyhf{}
\lhead{C++ Programming}
\rhead{Day 15}
\cfoot{\thepage}

\titleformat{\section}{\Large\bfseries\color{titleblue}}{\thesection.}{0.6em}{}
\titleformat{\subsection}{\large\bfseries}{\thesubsection.}{0.5em}{}

\lstdefinestyle{cppstyle}{
    language=C++,
    backgroundcolor=\color{codebg},
    frame=single,
    rulecolor=\color{codeframe},
    basicstyle=\ttfamily\small,
    keywordstyle=\color{keywordcolor}\bfseries,
    commentstyle=\color{commentcolor},
    stringstyle=\color{stringcolor},
    numbers=left,
    numberstyle=\tiny\color{gray},
    numbersep=8pt,
    tabsize=4,
    showstringspaces=false,
    breaklines=true,
    columns=fullflexible
}
\lstset{style=cppstyle}

\newtcolorbox{learningbox}{colback=blue!4,colframe=titleblue,title=\textbf{학습 목표},fonttitle=\bfseries}
\newtcolorbox{importantbox}{colback=yellow!8,colframe=orange!70!black,title=\textbf{핵심 포인트},fonttitle=\bfseries}
\newtcolorbox{practicebox}{colback=red!3,colframe=red!55!black,title=\textbf{프로젝트 목표},fonttitle=\bfseries}

\begin{document}

\begin{titlepage}
    \centering
    \vspace*{3cm}
    {\Huge\bfseries C++ Programming\par}
    \vspace{0.8cm}
    {\LARGE Day 15: Two-Body Gravity Simulator\par}
    \vspace{2cm}
    {\large OOP, Struct, Numerical Integration, Multi-file Project, CSV Output\par}
    \vfill
\end{titlepage}

\tableofcontents
\newpage

\section{학습 목표}
\begin{learningbox}
이 장을 마치면 다음 내용을 설명하고 직접 코드로 작성할 수 있다.
\begin{itemize}[leftmargin=1.5em]
    \item 두 입자의 상태를 \texttt{Particle} 클래스로 표현하기
    \item 위치, 속도, 가속도를 \texttt{struct}로 묶어 물리량을 명확하게 표현하기
    \item 두 입자의 상대위치로 중력가속도를 계산하기
    \item 두 입자의 운동방정식을 직접 수치적분하기
    \item Explicit Euler method로 한 timestep을 구현하기
    \item 헤더와 구현 파일을 분리한 다중 파일 프로젝트 구성하기
    \item 시뮬레이션 결과를 CSV 파일로 저장하기
    \item CSV 데이터를 Excel의 XY Scatter chart로 확인하기
\end{itemize}
\end{learningbox}

\section{프로젝트 개요}
Day 15는 C++ 문법을 추가로 학습하는 날이 아니라, 지금까지 배운 내용을 실제 물리 문제에 연결하는 미니프로젝트이다.
프로젝트의 목표는 두 질점 사이에 뉴턴 중력만 작용한다고 가정하고, 두 입자의 운동을 시간에 따라 수치적으로 계산하는 것이다.

\begin{practicebox}
두 입자의 질량, 위치, 속도를 초기조건으로 설정하고, 각 timestep에서 중력가속도를 계산한 뒤 Explicit Euler method로 위치와 속도를 갱신한다. 계산된 궤적은 \texttt{trajectory.csv}에 저장한다.
\end{practicebox}

\subsection{프로젝트 파일 구조}
\begin{verbatim}
day15/
|-- main.cpp
|-- Particle.h
|-- Particle.cpp
|-- Physics.h
|-- Physics.cpp
|-- Simulation.h
|-- Simulation.cpp
`-- output/
    `-- trajectory.csv
\end{verbatim}

\begin{importantbox}
이 프로젝트에서는 환산질량을 이용해 two-body 문제를 one-body 문제로 환원하지 않는다. 수치해석에서는 두 입자의 운동방정식을 각각 직접 적분할 수 있으므로, 두 \texttt{Particle} 객체의 상태를 그대로 갱신하는 구조를 사용한다.
\end{importantbox}

\section{물리 모델}
두 입자의 위치를 각각 $\mathbf r_1$, $\mathbf r_2$라 하고 상대위치를
\[
\mathbf r = \mathbf r_2 - \mathbf r_1
\]
로 정의한다.

거리 $r$은
\[
r = \sqrt{x^2+y^2}
\]
이다.

뉴턴의 만유인력 법칙에 따라 두 입자 사이 힘의 크기는
\[
F = G\frac{m_1m_2}{r^2}
\]
이고, 벡터형으로 쓰면
\[
\mathbf F_1 = Gm_1m_2\frac{\mathbf r}{r^3},
\qquad
\mathbf F_2 = -Gm_1m_2\frac{\mathbf r}{r^3}
\]
이다.

각 입자의 가속도는
\[
\mathbf a_1 = Gm_2\frac{\mathbf r}{r^3},
\qquad
\mathbf a_2 = -Gm_1\frac{\mathbf r}{r^3}
\]
이 된다.

\begin{importantbox}
$1/r^2$는 힘의 크기이고, 벡터식에서 $\mathbf r/r^3$가 나타나는 이유는 방향을 나타내는 단위벡터 $\hat{\mathbf r}=\mathbf r/r$를 곱했기 때문이다.
\end{importantbox}

\section{Particle 클래스}
\texttt{Particle}은 하나의 물체가 가진 질량, 위치, 속도를 저장한다.

\subsection{Particle.h}
\begin{lstlisting}
#ifndef PARTICLE_H
#define PARTICLE_H

class Particle
{
public:
    double mass;
    double x;
    double y;
    double vx;
    double vy;

    Particle(double mass, double x, double y, double vx, double vy);

    void showInfo() const;
};

#endif
\end{lstlisting}

\subsection{Particle.cpp}
\begin{lstlisting}
#include <iostream>
#include "Particle.h"

using namespace std;

Particle::Particle(double mass, double x, double y, double vx, double vy)
{
    this->mass = mass;
    this->x = x;
    this->y = y;
    this->vx = vx;
    this->vy = vy;
}

void Particle::showInfo() const
{
    cout << "Mass: " << mass << endl;
    cout << "Position: (" << x << ", " << y << ")" << endl;
    cout << "Velocity: (" << vx << ", " << vy << ")" << endl;
}
\end{lstlisting}

\section{Physics 모듈}
물리 계산은 \texttt{Particle} 클래스와 분리하여 \texttt{Physics.h}, \texttt{Physics.cpp}에서 처리한다.

\subsection{물리량 struct}
위치, 속도, 가속도는 서로 다른 의미를 가진 물리량이므로 별도의 구조체로 분리한다.

\begin{lstlisting}
struct Position
{
    double x;
    double y;
};

struct Velocity
{
    double vx;
    double vy;
};

struct Acceleration
{
    double ax;
    double ay;
};

struct Accelerations
{
    Acceleration acceleration1;
    Acceleration acceleration2;
};
\end{lstlisting}

\subsection{Physics.h}
\begin{lstlisting}
#ifndef PHYSICS_H
#define PHYSICS_H

#include "Particle.h"

struct Position
{
    double x;
    double y;
};

struct Velocity
{
    double vx;
    double vy;
};

struct Acceleration
{
    double ax;
    double ay;
};

struct Accelerations
{
    Acceleration acceleration1;
    Acceleration acceleration2;
};

Position getRelativePosition(
    const Particle& particle1,
    const Particle& particle2
);

Accelerations getAccelerations(
    const Particle& particle1,
    const Particle& particle2
);

#endif
\end{lstlisting}

\subsection{Physics.cpp}
\begin{lstlisting}
#include "Physics.h"
#include <cmath>

Position getRelativePosition(
    const Particle& particle1,
    const Particle& particle2
)
{
    Position relativePosition;
    relativePosition.x = particle2.x - particle1.x;
    relativePosition.y = particle2.y - particle1.y;

    return relativePosition;
}

Accelerations getAccelerations(
    const Particle& particle1,
    const Particle& particle2
)
{
    Position relativePosition = getRelativePosition(particle1, particle2);

    double x = relativePosition.x;
    double y = relativePosition.y;
    double m1 = particle1.mass;
    double m2 = particle2.mass;

    double r = std::sqrt(x * x + y * y);

    const double G = 6.67430e-11;

    Accelerations accelerations;

    accelerations.acceleration1.ax = G * m2 * x / (r * r * r);
    accelerations.acceleration1.ay = G * m2 * y / (r * r * r);

    accelerations.acceleration2.ax = -G * m1 * x / (r * r * r);
    accelerations.acceleration2.ay = -G * m1 * y / (r * r * r);

    return accelerations;
}
\end{lstlisting}

\begin{importantbox}
두 입자의 위치가 완전히 같아지면 $r=0$이 되어 분모가 0이 된다. 현재 프로젝트에서는 초기조건을 서로 다른 위치로 설정하고 충돌 처리는 구현하지 않는다.
\end{importantbox}

\section{Explicit Euler Method}
현재 상태에서 가속도를 계산한 뒤 timestep $\Delta t$만큼 전진한다.

위치는
\[
x_{n+1}=x_n+v_{x,n}\Delta t,
\qquad
y_{n+1}=y_n+v_{y,n}\Delta t
\]
로 갱신한다.

속도는
\[
v_{x,n+1}=v_{x,n}+a_{x,n}\Delta t,
\qquad
v_{y,n+1}=v_{y,n}+a_{y,n}\Delta t
\]
로 갱신한다.

\begin{importantbox}
Explicit Euler에서는 새 위치와 새 속도를 모두 기존 상태 $n$을 이용해 계산해야 한다. 속도를 먼저 갱신한 뒤 그 새 속도로 위치를 갱신하면 Semi-Implicit Euler가 된다.
\end{importantbox}

\section{Simulation 모듈}
\texttt{step()}은 두 입자의 상태를 reference로 받아 한 timestep만큼 직접 변경한다.

\subsection{Simulation.h}
\begin{lstlisting}
#ifndef SIMULATION_H
#define SIMULATION_H

#include "Particle.h"
#include "Physics.h"

void step(Particle& particle1, Particle& particle2, double dt);

#endif
\end{lstlisting}

\subsection{Simulation.cpp}
\begin{lstlisting}
#include "Simulation.h"

void step(Particle& particle1, Particle& particle2, double dt)
{
    Accelerations accelerations
        = getAccelerations(particle1, particle2);

    Position oldPosition1;
    oldPosition1.x = particle1.x;
    oldPosition1.y = particle1.y;

    Position oldPosition2;
    oldPosition2.x = particle2.x;
    oldPosition2.y = particle2.y;

    Velocity oldVelocity1;
    oldVelocity1.vx = particle1.vx;
    oldVelocity1.vy = particle1.vy;

    Velocity oldVelocity2;
    oldVelocity2.vx = particle2.vx;
    oldVelocity2.vy = particle2.vy;

    Position newPosition1;
    newPosition1.x = oldPosition1.x + oldVelocity1.vx * dt;
    newPosition1.y = oldPosition1.y + oldVelocity1.vy * dt;

    Position newPosition2;
    newPosition2.x = oldPosition2.x + oldVelocity2.vx * dt;
    newPosition2.y = oldPosition2.y + oldVelocity2.vy * dt;

    Velocity newVelocity1;
    newVelocity1.vx = oldVelocity1.vx
                    + accelerations.acceleration1.ax * dt;
    newVelocity1.vy = oldVelocity1.vy
                    + accelerations.acceleration1.ay * dt;

    Velocity newVelocity2;
    newVelocity2.vx = oldVelocity2.vx
                    + accelerations.acceleration2.ax * dt;
    newVelocity2.vy = oldVelocity2.vy
                    + accelerations.acceleration2.ay * dt;

    particle1.x = newPosition1.x;
    particle1.y = newPosition1.y;
    particle1.vx = newVelocity1.vx;
    particle1.vy = newVelocity1.vy;

    particle2.x = newPosition2.x;
    particle2.y = newPosition2.y;
    particle2.vx = newVelocity2.vx;
    particle2.vy = newVelocity2.vy;
}
\end{lstlisting}

\section{초기조건과 Simulation Loop}
테스트에서는 태양과 지구에 가까운 질량, 거리, 속도를 사용한다.

\begin{center}
\begin{tabular}{lll}
\toprule
항목 & Particle 1 & Particle 2 \\
\midrule
질량 & $1.989\times10^{30}$ kg & $5.972\times10^{24}$ kg \\
초기 위치 & $(0,0)$ m & $(1.496\times10^{11},0)$ m \\
초기 속도 & $(0,0)$ m/s & $(0,29780)$ m/s \\
\bottomrule
\end{tabular}
\end{center}

시간 간격은
\[
\Delta t = 3600\ \text{s}
\]
로 두어 한 번의 \texttt{step()}이 1시간을 의미하도록 한다.

\section{CSV 출력과 main.cpp}
CSV 파일의 첫 행은 각 열의 의미를 나타낸다.

\begin{verbatim}
time,x1,y1,x2,y2
\end{verbatim}

\subsection{main.cpp}
\begin{lstlisting}
#include <iostream>
#include <fstream>
#include "Particle.h"
#include "Physics.h"
#include "Simulation.h"

using namespace std;

int main()
{
    Particle particle1(
        1.989e30,
        0.0, 0.0,
        0.0, 0.0
    );

    Particle particle2(
        5.972e24,
        1.496e11, 0.0,
        0.0, 29780.0
    );

    double dt = 3600.0;
    double totalTime = 365.25 * 24.0 * 3600.0;

    int steps = static_cast<int>(totalTime / dt);

    ofstream file("output/trajectory.csv");

    if (!file.is_open())
    {
        cout << "Failed to open file" << endl;
        return 1;
    }

    file << "time,x1,y1,x2,y2" << endl;

    for (int i = 0; i <= steps; i++)
    {
        double time = i * dt;

        file << time << ","
             << particle1.x << ","
             << particle1.y << ","
             << particle2.x << ","
             << particle2.y << endl;

        if (i < steps)
        {
            step(particle1, particle2, dt);
        }
    }

    file.close();

    cout << "Simulation completed." << endl;
    cout << "Data saved to output/trajectory.csv" << endl;

    return 0;
}
\end{lstlisting}

\section{Build와 실행}
모든 구현 파일을 함께 컴파일해야 한다.

\begin{lstlisting}
g++ main.cpp Particle.cpp Physics.cpp Simulation.cpp -o simulation
\end{lstlisting}

실행 후 \texttt{output/trajectory.csv}가 생성되면 시뮬레이션이 정상적으로 완료된 것이다.

\begin{importantbox}
\texttt{main.cpp}만 컴파일하면 \texttt{Particle.cpp}, \texttt{Physics.cpp}, \texttt{Simulation.cpp}에 있는 구현을 linker가 찾지 못해 \texttt{undefined reference} 오류가 발생한다.
\end{importantbox}

\section{CSV 데이터로 궤도 확인}
CSV 자체는 그래프를 저장하는 형식이 아니므로 Excel에서 데이터를 불러온 뒤 XY Scatter chart를 사용한다.

Particle 2의 궤도는
\begin{itemize}[leftmargin=1.5em]
    \item X 값: \texttt{x2}
    \item Y 값: \texttt{y2}
\end{itemize}
로 설정한다.

\begin{importantbox}
Excel이 과학적 표기법 데이터를 문자열로 인식하면 셀 앞에 작은따옴표가 붙을 수 있다. 이 경우 숫자로 변환한 뒤 XY Scatter chart를 생성해야 한다. 단순 Line chart를 사용하면 행 번호나 시간에 대한 그래프가 되어 실제 궤도가 아니다.
\end{importantbox}

\section{프로젝트에서 사용한 C++ 개념}
\begin{center}
\begin{tabular}{ll}
\toprule
개념 & 프로젝트에서의 역할 \\
\midrule
Class & Particle 객체 정의 \\
Struct & Position, Velocity, Acceleration 데이터 묶음 \\
Reference & \texttt{step()}에서 원본 Particle 상태 변경 \\
Const reference & 물리 계산 함수에서 불필요한 복사 방지 \\
Header / cpp 분리 & 선언과 구현 분리 \\
Function & 물리 계산과 시뮬레이션 역할 분리 \\
File I/O & CSV 출력 \\
Loop & timestep 반복 \\
\texttt{static\_cast} & 반복 횟수를 정수로 변환 \\
\bottomrule
\end{tabular}
\end{center}

\section{Numerical Method의 한계}
Explicit Euler method는 구현이 단순하지만 장시간 궤도 계산에서 에너지 오차가 누적될 수 있다.
따라서 원에 가까운 초기조건을 사용해도 궤도가 완전히 닫히지 않거나 점점 변형될 수 있다.

향후 확장에서는 다음 방법을 비교할 수 있다.
\begin{itemize}[leftmargin=1.5em]
    \item timestep $\Delta t$ 감소
    \item Semi-Implicit Euler
    \item Velocity Verlet
    \item Runge--Kutta method
    \item total energy와 angular momentum의 시간 변화 확인
\end{itemize}

\section{Day 15 정리}
Day 15에서는 두 입자 중력 문제를 C++ 다중 파일 프로젝트로 구현했다.
핵심은 물리식을 코드로 직접 변환하고, 각 timestep에서 두 입자의 가속도, 속도, 위치를 순서대로 계산하는 것이다.

프로그램의 전체 흐름은 다음과 같다.

\begin{verbatim}
Initial Conditions
       |
       v
Particle 1 + Particle 2
       |
       v
Relative Position
       |
       v
Gravitational Accelerations
       |
       v
Explicit Euler step
       |
       v
Updated Particle States
       |
       v
CSV Output
       |
       v
Repeat
\end{verbatim}

이 프로젝트를 통해 클래스, 구조체, 참조, 함수 분리, 헤더 파일, 수치해석, 파일 입출력을 하나의 물리 시뮬레이션 프로그램 안에서 연결했다.

\end{document}
